\documentclass[10pt,a4paper]{amsart}
\usepackage{amsmath,amssymb,amsthm,graphicx}
\usepackage{verbatim}
\newenvironment{reference}{\comment}{\endcomment}
%\newenvironment{reference}{}{}
\newenvironment{slogan}{\comment}{\endcomment}
\newenvironment{history}{\comment}{\endcomment}
% For commutative diagrams we use Xy-pic
\usepackage[all]{xy}

% We use 2cell for 2-commutative diagrams.
\xyoption{2cell}
% Theorem environments.
%
\theoremstyle{plain}
\newtheorem{theorem}[subsection]{Theorem}
\newtheorem{proposition}[subsection]{Proposition}
\newtheorem{lemma}[subsection]{Lemma}

\theoremstyle{definition}
\newtheorem{definition}[subsection]{Definition}
\newtheorem{example}[subsection]{Example}
\newtheorem{exercise}[subsection]{Exercise}
\newtheorem{situation}[subsection]{Situation}

\theoremstyle{remark}
\newtheorem{remark}[subsection]{Remark}
\newtheorem{remarks}[subsection]{Remarks}

\numberwithin{equation}{subsection}

% Macros
%
\def\lim{\mathop{\mathrm{lim}}\nolimits}
\def\colim{\mathop{\mathrm{colim}}\nolimits}
\def\Spec{\mathop{\mathrm{Spec}}}
\def\Hom{\mathop{\mathrm{Hom}}\nolimits}
\def\Ext{\mathop{\mathrm{Ext}}\nolimits}
\def\SheafHom{\mathop{\mathcal{H}\!\mathit{om}}\nolimits}
\def\SheafExt{\mathop{\mathcal{E}\!\mathit{xt}}\nolimits}
\def\Sch{\mathit{Sch}}
\def\Mor{\mathop{\mathrm{Mor}}\nolimits}
\def\Ob{\mathop{\mathrm{Ob}}\nolimits}
\def\Sh{\mathop{\mathit{Sh}}\nolimits}
\def\NL{\mathop{N\!L}\nolimits}
\def\CH{\mathop{\mathrm{CH}}\nolimits}
\def\proetale{{pro\text{-}\acute{e}tale}}
\def\etale{{\acute{e}tale}}
\def\QCoh{\mathit{QCoh}}
\def\Ker{\mathop{\mathrm{Ker}}}
\def\Im{\mathop{\mathrm{Im}}}
\def\Coker{\mathop{\mathrm{Coker}}}
\def\Coim{\mathop{\mathrm{Coim}}}

% Boxtimes
%
\DeclareMathSymbol{\boxtimes}{\mathbin}{AMSa}{"02}

%
% Macros for moduli stacks/spaces
%
\def\QCohstack{\mathcal{QC}\!\mathit{oh}}
\def\Cohstack{\mathcal{C}\!\mathit{oh}}
\def\Spacesstack{\mathcal{S}\!\mathit{paces}}
\def\Quotfunctor{\mathrm{Quot}}
\def\Hilbfunctor{\mathrm{Hilb}}
\def\Curvesstack{\mathcal{C}\!\mathit{urves}}
\def\Polarizedstack{\mathcal{P}\!\mathit{olarized}}
\def\Complexesstack{\mathcal{C}\!\mathit{omplexes}}
% \Pic is the operator that assigns to X its picard group, usage \Pic(X)
% \Picardstack_{X/B} denotes the Picard stack of X over B
% \Picardfunctor_{X/B} denotes the Picard functor of X over B
\def\Pic{\mathop{\mathrm{Pic}}\nolimits}
\def\Picardstack{\mathcal{P}\!\mathit{ic}}
\def\Picardfunctor{\mathrm{Pic}}
\def\Deformationcategory{\mathcal{D}\!\mathit{ef}}
\usepackage{cleveref}
\crefname{originalreference}{original reference}{original references}
\Crefname{originalreference}{Original reference}{Original references}
\makeatletter
\newlabel{algebra-lemma-flat-going-down}{{algebra-lemma-flat-going-down}{source}{}{sourcepointer-1}{}}
\newlabel{algebra-lemma-flat-going-down@cref}{{[originalreference][0][]algebra-lemma-flat-going-down}{[1][0][]source}{}{}{}}
\makeatother
\crefname{originallemma}{original lemma}{original lemmas}
\Crefname{originallemma}{Original lemma}{Original lemmas}
\makeatletter
\newlabel{lemma-apply-CM}{{lemma-apply-CM}{source}{}{sourcepointer-2}{}}
\newlabel{lemma-apply-CM@cref}{{[originallemma][0][]lemma-apply-CM}{[1][0][]source}{}{}{}}
\makeatother
\crefname{originallemma}{original lemma}{original lemmas}
\Crefname{originallemma}{Original lemma}{Original lemmas}
\makeatletter
\newlabel{lemma-base-change-relative-algebraic}{{lemma-base-change-relative-algebraic}{source}{}{sourcepointer-3}{}}
\newlabel{lemma-base-change-relative-algebraic@cref}{{[originallemma][0][]lemma-base-change-relative-algebraic}{[1][0][]source}{}{}{}}
\makeatother
\crefname{originallemma}{original lemma}{original lemmas}
\Crefname{originallemma}{Original lemma}{Original lemmas}
\makeatletter
\newlabel{lemma-depth-dualizing-module}{{lemma-depth-dualizing-module}{source}{}{sourcepointer-4}{}}
\newlabel{lemma-depth-dualizing-module@cref}{{[originallemma][0][]lemma-depth-dualizing-module}{[1][0][]source}{}{}{}}
\makeatother
\crefname{originallemma}{original lemma}{original lemmas}
\Crefname{originallemma}{Original lemma}{Original lemmas}
\makeatletter
\newlabel{lemma-relative-dualizing-if-have-omega}{{lemma-relative-dualizing-if-have-omega}{source}{}{sourcepointer-5}{}}
\newlabel{lemma-relative-dualizing-if-have-omega@cref}{{[originallemma][0][]lemma-relative-dualizing-if-have-omega}{[1][0][]source}{}{}{}}
\makeatother
\crefname{originallemma}{original lemma}{original lemmas}
\Crefname{originallemma}{Original lemma}{Original lemmas}
\makeatletter
\newlabel{lemma-shriek-normalized}{{lemma-shriek-normalized}{source}{}{sourcepointer-6}{}}
\newlabel{lemma-shriek-normalized@cref}{{[originallemma][0][]lemma-shriek-normalized}{[1][0][]source}{}{}{}}
\makeatother
\crefname{originallemma}{original lemma}{original lemmas}
\Crefname{originallemma}{Original lemma}{Original lemmas}
\makeatletter
\newlabel{lemma-sitting-in-degrees}{{lemma-sitting-in-degrees}{source}{}{sourcepointer-7}{}}
\newlabel{lemma-sitting-in-degrees@cref}{{[originallemma][0][]lemma-sitting-in-degrees}{[1][0][]source}{}{}{}}
\makeatother
\crefname{originallemma}{original lemma}{original lemmas}
\Crefname{originallemma}{Original lemma}{Original lemmas}
\makeatletter
\newlabel{lemma-ubiquity-dualizing}{{lemma-ubiquity-dualizing}{source}{}{sourcepointer-8}{}}
\newlabel{lemma-ubiquity-dualizing@cref}{{[originallemma][0][]lemma-ubiquity-dualizing}{[1][0][]source}{}{}{}}
\makeatother
\crefname{originalreference}{original reference}{original references}
\Crefname{originalreference}{Original reference}{Original references}
\makeatletter
\newlabel{more-algebra-lemma-completion-CM}{{more-algebra-lemma-completion-CM}{source}{}{sourcepointer-9}{}}
\newlabel{more-algebra-lemma-completion-CM@cref}{{[originalreference][0][]more-algebra-lemma-completion-CM}{[1][0][]source}{}{}{}}
\makeatother
\crefname{originalreference}{original reference}{original references}
\Crefname{originalreference}{Original reference}{Original references}
\makeatletter
\newlabel{more-algebra-lemma-hom-into-S2}{{more-algebra-lemma-hom-into-S2}{source}{}{sourcepointer-10}{}}
\newlabel{more-algebra-lemma-hom-into-S2@cref}{{[originalreference][0][]more-algebra-lemma-hom-into-S2}{[1][0][]source}{}{}{}}
\makeatother
\crefname{originalreference}{original reference}{original references}
\Crefname{originalreference}{Original reference}{Original references}
\makeatletter
\newlabel{morphisms-lemma-dimension-fibre-after-base-change}{{morphisms-lemma-dimension-fibre-after-base-change}{source}{}{sourcepointer-11}{}}
\newlabel{morphisms-lemma-dimension-fibre-after-base-change@cref}{{[originalreference][0][]morphisms-lemma-dimension-fibre-after-base-change}{[1][0][]source}{}{}{}}
\makeatother
\crefname{originalreference}{original reference}{original references}
\Crefname{originalreference}{Original reference}{Original references}
\makeatletter
\newlabel{morphisms-lemma-dimension-formula-general}{{morphisms-lemma-dimension-formula-general}{source}{}{sourcepointer-12}{}}
\newlabel{morphisms-lemma-dimension-formula-general@cref}{{[originalreference][0][]morphisms-lemma-dimension-formula-general}{[1][0][]source}{}{}{}}
\makeatother
\title{A vanishing range and an S2 statement for the relative dualizing complex over a Cohen-Macaulay base}
\author{The Stacks Project Authors}
\date{stacks-project:dualizing.tex:local-snapshot-20261012; GNU FDL 1.2 or later}
\begin{document}
\maketitle
\noindent\textit{Editorial note.} All delivered mathematical text is the author's own, copied from the pinned Stacks Project chapter source sources/stacks-dualizing/pinned-original-source.tex; no statement or proof step was rewritten, reordered or removed.\par\medskip
\noindent\textit{Editorial note.} Preamble provenance: the Stacks Project keeps one single shared preamble for every chapter; this card takes it directly from the staged sources/stacks-dualizing/preamble.tex as three verbatim spans (lines 8-12, 14-18 and 154-217), so no derived file was written and nothing was added or redefined. The wrapper class and the base packages are supplied by the assembler.\par\medskip
\noindent\textit{Editorial note.} Bibliography provenance: neither the counted statement nor its proof contains a citation (this lane rescans the whole span before writing the card), so the card carries no bibliography and no entry was reconstructed.\par\medskip
\noindent\textit{Editorial note.} Selection provenance: this card was opened under the file-based control contract (CONTROL.md), with the mechanical pre-card gate run on its statement and proof interval before the card was written.\par\medskip
\noindent\textit{Allowed standard prerequisites (editorial).}\begin{itemize}
\item The base change of the relative dualizing complex along a localization of the base. Invoked at the start of the counted proof and again in the middle (source lines 5366-5367 and 5383-5384). Reduces to a Cohen-Macaulay local base.
\item The Cohen-Macaulayness of the completion of a Cohen-Macaulay local ring. Invoked inside the counted proof (source lines 5373-5375). Allows replacing the base by its completion.
\item The existence of a dualizing complex after passing to a suitable local base. Invoked inside the counted proof (source lines 5385-5388). Makes a dualizing complex available for the base.
\item The normalized form of the dualizing complex of a Cohen-Macaulay local ring. Invoked inside the counted proof (source lines 5399-5402). Writes the base dualizing complex as a module in a single degree.
\item The degree confinement and the depth statement for a normalized dualizing complex. Invoked inside the counted proof (source lines 5416-5418). Gives the vanishing range of the cohomology of the shriek of the base dualizing complex.
\end{itemize}\medskip
\ifdefined\hypertarget\hypertarget{sourcepointer-1}{}\fi
\noindent\textit{Source reference.} Original reference algebra-lemma-flat-going-down: sources/stacks-dualizing/pinned-original-source.tex, source lines 5379--5379. Navigation pointer only: the counted text refers to the Stacks item algebra-lemma-flat-going-down, which belongs to another chapter of the project and is not part of this card. The pointer carries the target's own label as its number because this chapter's staged snapshot does not contain the other chapter's tag table; its content is not reproduced here.\par\medskip
\ifdefined\hypertarget\hypertarget{sourcepointer-2}{}\fi
\noindent\textit{Source reference.} Original lemma lemma-apply-CM: sources/stacks-dualizing/pinned-original-source.tex, source lines 3842--3842. Navigation pointer only: the counted text refers to the Stacks item lemma-apply-CM, which is an existing labeled result of this same chapter and is not part of this card. The pointer carries the target's own label as its number and its definition line as its span; its content is not reproduced here.\par\medskip
\ifdefined\hypertarget\hypertarget{sourcepointer-3}{}\fi
\noindent\textit{Source reference.} Original lemma lemma-base-change-relative-algebraic: sources/stacks-dualizing/pinned-original-source.tex, source lines 5023--5023. Navigation pointer only: the counted text refers to the Stacks item lemma-base-change-relative-algebraic, which is an existing labeled result of this same chapter and is not part of this card. The pointer carries the target's own label as its number and its definition line as its span; its content is not reproduced here.\par\medskip
\ifdefined\hypertarget\hypertarget{sourcepointer-4}{}\fi
\noindent\textit{Source reference.} Original lemma lemma-depth-dualizing-module: sources/stacks-dualizing/pinned-original-source.tex, source lines 3576--3576. Navigation pointer only: the counted text refers to the Stacks item lemma-depth-dualizing-module, which is an existing labeled result of this same chapter and is not part of this card. The pointer carries the target's own label as its number and its definition line as its span; its content is not reproduced here.\par\medskip
\ifdefined\hypertarget\hypertarget{sourcepointer-5}{}\fi
\noindent\textit{Source reference.} Original lemma lemma-relative-dualizing-if-have-omega: sources/stacks-dualizing/pinned-original-source.tex, source lines 4911--4911. Navigation pointer only: the counted text refers to the Stacks item lemma-relative-dualizing-if-have-omega, which is an existing labeled result of this same chapter and is not part of this card. The pointer carries the target's own label as its number and its definition line as its span; its content is not reproduced here.\par\medskip
\ifdefined\hypertarget\hypertarget{sourcepointer-6}{}\fi
\noindent\textit{Source reference.} Original lemma lemma-shriek-normalized: sources/stacks-dualizing/pinned-original-source.tex, source lines 5221--5221. Navigation pointer only: the counted text refers to the Stacks item lemma-shriek-normalized, which is an existing labeled result of this same chapter and is not part of this card. The pointer carries the target's own label as its number and its definition line as its span; its content is not reproduced here.\par\medskip
\ifdefined\hypertarget\hypertarget{sourcepointer-7}{}\fi
\noindent\textit{Source reference.} Original lemma lemma-sitting-in-degrees: sources/stacks-dualizing/pinned-original-source.tex, source lines 3232--3232. Navigation pointer only: the counted text refers to the Stacks item lemma-sitting-in-degrees, which is an existing labeled result of this same chapter and is not part of this card. The pointer carries the target's own label as its number and its definition line as its span; its content is not reproduced here.\par\medskip
\ifdefined\hypertarget\hypertarget{sourcepointer-8}{}\fi
\noindent\textit{Source reference.} Original lemma lemma-ubiquity-dualizing: sources/stacks-dualizing/pinned-original-source.tex, source lines 4343--4343. Navigation pointer only: the counted text refers to the Stacks item lemma-ubiquity-dualizing, which is an existing labeled result of this same chapter and is not part of this card. The pointer carries the target's own label as its number and its definition line as its span; its content is not reproduced here.\par\medskip
\ifdefined\hypertarget\hypertarget{sourcepointer-9}{}\fi
\noindent\textit{Source reference.} Original reference more-algebra-lemma-completion-CM: sources/stacks-dualizing/pinned-original-source.tex, source lines 5375--5375. Navigation pointer only: the counted text refers to the Stacks item more-algebra-lemma-completion-CM, which belongs to another chapter of the project and is not part of this card. The pointer carries the target's own label as its number because this chapter's staged snapshot does not contain the other chapter's tag table; its content is not reproduced here.\par\medskip
\ifdefined\hypertarget\hypertarget{sourcepointer-10}{}\fi
\noindent\textit{Source reference.} Original reference more-algebra-lemma-hom-into-S2: sources/stacks-dualizing/pinned-original-source.tex, source lines 5429--5429. Navigation pointer only: the counted text refers to the Stacks item more-algebra-lemma-hom-into-S2, which belongs to another chapter of the project and is not part of this card. The pointer carries the target's own label as its number because this chapter's staged snapshot does not contain the other chapter's tag table; its content is not reproduced here.\par\medskip
\ifdefined\hypertarget\hypertarget{sourcepointer-11}{}\fi
\noindent\textit{Source reference.} Original reference morphisms-lemma-dimension-fibre-after-base-change: sources/stacks-dualizing/pinned-original-source.tex, source lines 5382--5382. Navigation pointer only: the counted text refers to the Stacks item morphisms-lemma-dimension-fibre-after-base-change, which belongs to another chapter of the project and is not part of this card. The pointer carries the target's own label as its number because this chapter's staged snapshot does not contain the other chapter's tag table; its content is not reproduced here.\par\medskip
\ifdefined\hypertarget\hypertarget{sourcepointer-12}{}\fi
\noindent\textit{Source reference.} Original reference morphisms-lemma-dimension-formula-general: sources/stacks-dualizing/pinned-original-source.tex, source lines 5410--5410. Navigation pointer only: the counted text refers to the Stacks item morphisms-lemma-dimension-formula-general, which belongs to another chapter of the project and is not part of this card. The pointer carries the target's own label as its number because this chapter's staged snapshot does not contain the other chapter's tag table; its content is not reproduced here.\par\medskip
\begin{lemma}
\label{lemma-relative-dualizing-CM-vanishing}
Let $R \to A$ be a finite type homomorphism of Noetherian rings.
Let $\mathfrak p \subset R$ be a prime ideal. Assume
\begin{enumerate}
\item $R_\mathfrak p$ is Cohen-Macaulay, and
\item for any minimal prime $\mathfrak q \subset A$ we have
$\text{trdeg}_{\kappa(R \cap \mathfrak q)} \kappa(\mathfrak q) \leq r$.
\end{enumerate}
Then
$$
H^i(\omega_{A/R}^\bullet)_\mathfrak p \not = 0 \Rightarrow - r \leq i
$$
and $H^{-r}(\omega_{A/R}^\bullet)_\mathfrak p$ is $(S_2)$
as an $A_\mathfrak p$-module.
\end{lemma}

\begin{proof}
We may replace $R$ by $R_\mathfrak p$ by
Lemma \ref{lemma-base-change-relative-algebraic}.
Thus we may assume $R$ is a Cohen-Macaulay local ring
and we have to show the assertions of the lemma
for the $A$-modules $H^i(\omega_{A/R}^\bullet)$.

\medskip\noindent
Let $R^\wedge$ be the completion of $R$.
The map $R \to R^\wedge$ is flat and $R^\wedge$ is Cohen-Macaulay
(More on Algebra, Lemma \ref{more-algebra-lemma-completion-CM}).
Observe that the minimal primes of $A \otimes_R R^\wedge$
lie over minimal primes of $A$ by the flatness of
$A \to A \otimes_R R^\wedge$ (and going down for flatness, see
Algebra, Lemma \ref{algebra-lemma-flat-going-down}).
Thus condition (2) holds for the finite type ring map
$R^\wedge \to A \otimes_R R^\wedge$ by
Morphisms, Lemma \ref{morphisms-lemma-dimension-fibre-after-base-change}.
Appealing to Lemma \ref{lemma-base-change-relative-algebraic}
once again it suffices to prove the lemma for
$R^\wedge \to A \otimes_R R^\wedge$. In this way, using
Lemma \ref{lemma-ubiquity-dualizing},
we may assume $R$ is a Noetherian local
Cohen-Macaulay ring which has a dualizing complex $\omega_R^\bullet$.

\medskip\noindent
Let $\mathfrak m \subset A$ be a maximal ideal.
It suffices to show that the assertions of
the lemma hold for $H^i(\omega_{A/R}^\bullet)_\mathfrak m$.
If $\mathfrak m$ does not lie over the maximal ideal of $R$,
then we replace $R$ by a localization to reduce to this case
(small detail omitted). 

\medskip\noindent
We may assume $\omega_R^\bullet$ is normalized.
Setting $d = \dim(R)$ we see that $\omega_R^\bullet = \omega_R[d]$
for some $R$-module $\omega_R$, see
Lemma \ref{lemma-apply-CM}. Set
$\omega_A^\bullet = \varphi^!(\omega_R^\bullet)$.
By Lemma \ref{lemma-relative-dualizing-if-have-omega} we have
$$
\omega_{A/R}^\bullet =
R\Hom_A(\omega_R[d] \otimes_R^\mathbf{L} A, \omega_A^\bullet)
$$
By the dimension formula we have $\dim(A_\mathfrak m) \leq d + r$, see
Morphisms, Lemma \ref{morphisms-lemma-dimension-formula-general}
and use that $\kappa(\mathfrak m)$ is finite over the residue field of $R$
by the Hilbert Nullstellensatz.
By Lemma \ref{lemma-shriek-normalized}
we see that $(\omega_A^\bullet)_\mathfrak m$
is a normalized dualizing complex for $A_\mathfrak m$.
Hence $H^i((\omega_A^\bullet)_\mathfrak m)$ is nonzero
only for $-d - r \leq i \leq 0$, see
Lemma \ref{lemma-sitting-in-degrees}.
Since $\omega_R[d] \otimes_R^\mathbf{L} A$ lives in
degrees $\leq -d$ we conclude the vanishing holds.
Finally, we also see that
$$
H^{-r}(\omega_{A/R}^\bullet)_\mathfrak m =
\Hom_A(\omega_R \otimes_R A, H^{-d - r}(\omega_A^\bullet))_\mathfrak m
$$
Since $H^{-d - r}(\omega_A^\bullet)_\mathfrak m$ is $(S_2)$ by
Lemma \ref{lemma-depth-dualizing-module}
we find that the final statement is true by
More on Algebra, Lemma \ref{more-algebra-lemma-hom-into-S2}.
\end{proof}

\end{document}
